\SourceLeaf{234}%
das Krümmungsmaß
\[
k=\frac{\mathrm{d}X\mathord{.}\delta Y-\mathrm{d}Y\mathord{.}\delta X}{\mathrm{d}x\mathord{.}\delta y-\mathrm{d}y\mathord{.}\delta x}
\]
Nehmen wir nun an, die gekrümmte Fläche sei nach der dritten in Artikel 4 behandelten Methode gegeben, so sind \(X\) und \(Y\) Funktionen der Größen \(x,y\). Daraus ergibt sich
\[
\begin{aligned}
\mathrm{d}X&=(\tfrac{\mathrm{d}X}{\mathrm{d}x})\mathrm{d}x+(\tfrac{\mathrm{d}X}{\mathrm{d}y})\mathrm{d}y\\
\delta X&=(\tfrac{\mathrm{d}X}{\mathrm{d}x})\delta x+(\tfrac{\mathrm{d}X}{\mathrm{d}y})\delta y\\
\mathrm{d}Y&=(\tfrac{\mathrm{d}Y}{\mathrm{d}x})\mathrm{d}x+(\tfrac{\mathrm{d}Y}{\mathrm{d}y})\mathrm{d}y\\
\delta Y&=(\tfrac{\mathrm{d}Y}{\mathrm{d}x})\delta x+(\tfrac{\mathrm{d}Y}{\mathrm{d}y})\delta y
\end{aligned}
\]
Setzt man diese Werte in den vorhergehenden Ausdruck ein, so erhält man
\[
k=(\tfrac{\mathrm{d}X}{\mathrm{d}x})(\tfrac{\mathrm{d}Y}{\mathrm{d}y})-(\tfrac{\mathrm{d}X}{\mathrm{d}y})(\tfrac{\mathrm{d}Y}{\mathrm{d}x})
\]
\par\indent Setzen wir wie oben
\[
\frac{\mathrm{d}z}{\mathrm{d}x}=t,\qquad\frac{\mathrm{d}z}{\mathrm{d}y}=u
\]
und außerdem
\[
\frac{\mathrm{d}\mathrm{d}z}{\mathrm{d}x^2}=T,\qquad\frac{\mathrm{d}\mathrm{d}z}{\mathrm{d}x\mathord{.}\mathrm{d}y}=U,\qquad\frac{\mathrm{d}\mathrm{d}z}{\mathrm{d}y^2}=V
\]
oder
\[
\mathrm{d}t=T\mathrm{d}x+U\mathrm{d}y,\qquad\mathrm{d}u=U\mathrm{d}x+V\mathrm{d}y
\]
so erhalten wir aus den oben angegebenen Formeln
\[
X=-tZ,\qquad Y=-uZ,\qquad(1+tt+uu)ZZ=1
\]
und daraus
\[
\begin{aligned}
&\mathrm{d}X=-Z\mathrm{d}t-t\mathrm{d}Z\\
&\mathrm{d}Y=-Z\mathrm{d}u-u\mathrm{d}Z\\
&(1+tt+uu)\mathrm{d}Z+Z(t\mathrm{d}t+u\mathrm{d}u)=0
\end{aligned}
\]
oder
\[
\begin{aligned}
\mathrm{d}Z&=-Z^3(t\mathrm{d}t+u\mathrm{d}u)\\
\mathrm{d}X&=-Z^3(1+uu)\mathrm{d}t+Z^3tu\mathrm{d}u\\
\mathrm{d}Y&=+Z^3tu\mathrm{d}t-Z^3(1+tt)\mathrm{d}u\mathord{.}
\end{aligned}
\]
